\clearpage
\chapter*{Erratum and supplement to Exposé no. 236}
\addcontentsline{toc}{chapter}{Erratum and supplement to Exposé no. 236}
\src{303}

\noindent\textbf{8. TDTE VI: Picard schemes. General properties.}

\begin{center}
[Bourbaki Seminar, vol. 14, 1961/62, no. 236, 23 pp.]
\end{center}

\noindent\textbf{Page 12, line 16.} --- For “affine line,” read “affine line
with the origin removed.”

\noindent\textbf{Line 18.} --- For “$X[t]$,” read “$X[t,t^{-1}]$.”

\noindent\textbf{Page 16, Proposition 3.1, sixth line.} --- For the malformed
construction “The hypothesis ... results,” read “It follows from the
hypothesis ....”

\noindent\textbf{Page 21, §4: The finiteness theorem for the Picard scheme.}
--- Finiteness questions of the kind raised in this section have been almost
completely resolved since the present exposé was written. We indicate the main
facts now known in this direction. To simplify the statements, we implicitly
assume that all the Picard preschemes occurring in them exist, although an
evident modification of the statements removes every existence hypothesis.
In what follows, $S$ denotes a noetherian scheme and $X,Y$ proper schemes over
$S$.

\medskip
\noindent(i) Let $f:X\longrightarrow Y$ be a surjective morphism. Then
\[
f^*: \operatorname{Pic}_{Y/S}\longrightarrow\operatorname{Pic}_{X/S}
\]
is a morphism of finite type.

\src{304}
\noindent(ii) Suppose that $Y$ is projective over $S$ and carries an invertible
Module relatively ample over $S$. Let $X$ be the zero scheme of an arbitrary
section of this Module, and let $f:X\longrightarrow Y$ be the canonical
immersion. Finally, suppose that the irreducible components of the fibers of
$Y/S$ have dimension $\geq3$. Then
\[
f^*: \operatorname{Pic}_{Y/S}\longrightarrow\operatorname{Pic}_{X/S}
\]
is of finite type.

\noindent(iii) Suppose that $X$ is projective over $S$ and all its geometric
fibers are integral of dimension $n$. Let $\mathcal O_X(1)$ be an ample
invertible Module on $X$, permitting Hilbert polynomials to be defined. If $M$
is a subset of $\operatorname{Pic}_{X/S}$, then $M$ is quasi-compact if and
only if, in the Hilbert polynomials
\[
a_0x^n+a_1x^{n-1}+a_2x^{n-2}+\cdots
\]
of the elements of $M$, the coefficients $a_1$ and $a_2$ remain bounded, omitting every coefficient whose monomial would have negative degree; thus only $a_1$ occurs when $n=1$, and neither coefficient occurs when $n=0$.

\noindent(iv) For every integer $n\neq0$, the $n$th-power homomorphism in the
group prescheme $\operatorname{Pic}_{X/S}$ is a morphism of finite type.

Notice that (i) and (ii) also say that, under their respective hypotheses, a
subset $M$ of $\operatorname{Pic}_{Y/S}$ is quasi-compact if and only if its
image in $\operatorname{Pic}_{X/S}$ is quasi-compact. It follows that an
invertible Module $\mathcal L$ on $Y$ is $\tau$-equivalent to zero if and only
if its inverse image on $X$ is; in other words,
$\operatorname{Pic}_{Y/S}^\tau$ is the inverse image of
$\operatorname{Pic}_{X/S}^\tau$. In particular, to show that the former
prescheme is of finite type over $S$, it is enough to prove this for the latter,
since $\operatorname{Pic}_{Y/S}\longrightarrow
\operatorname{Pic}_{X/S}$ is of finite type. Applying (i), Chow's lemma, and
(iii), one obtains:

\medskip
\noindent(v) $\operatorname{Pic}_{X/S}^\tau$ is of finite type over $S$.

In general, combining (i) for a finite morphism with (ii) reduces most
finiteness questions to the case in which $X/S$ has geometrically integral
normal fibers of dimension $\leq2$. Often, by applying (i) to a surjective
morphism not necessarily finite and using resolution of singularities of
algebraic surfaces (due in arbitrary characteristic to Abhyankar), one even
reduces to the case in which $X/S$ is smooth, and hence has nonsingular
geometric fibers of dimension 2. Together with (v) and the Picard--Igusa
inequality bounding the rank of the Néron--Severi group of a nonsingular
projective surface, this proves the following generalization of Néron's
finiteness theorem:

\medskip
\noindent(vi) Let $X/S$ be proper over $S$ and otherwise arbitrary. Then the
Néron--Severi groups
\[
\operatorname{Pic}_{X_i/k_i}/\operatorname{Pic}_{X_i/k_i}^\circ
\]
of the geometric fibers $X_i/k_i$ of $X/S$

\src{305}
are of finite type, and their ranks and the orders of their torsion subgroups
remain bounded.

The same reduction to nonsingular surfaces, together with the theorems known
in that case (namely Néron's finiteness theorem and the nondegeneracy of the
intersection form on the Néron--Severi group), implies:

\medskip
\noindent(vii) Let $X$ be a proper scheme over an algebraically closed field.
There are finitely many closed integral curves $C_i$ ($1\leq i\leq r$) in
$X$ with the following property: for every subset $M$ of
$\operatorname{Pic}_{X/k}$, $M$ is quasi-compact if and only if the integers
\[
\deg_{C_i'}\mathcal L\qquad(\mathcal L\in M)
\]
remain bounded, where $C_i'$ denotes the normalization of $C_i$.

Here $r$ may be taken to be the rank of the Néron--Severi group. Once that
group is known to be of finite type, (vii) reduces to the fact that the linear
forms on the Néron--Severi group defined by curves $C$ in $X$ vanish
simultaneously only on its torsion elements. When $X$ is projective and
nonsingular, this result, as well as (v), was due to Matsusaka. Using (vii),
one readily obtains the following characterization of invertible Modules
$\tau$-equivalent to zero on $X$:

\medskip
\noindent(viii) Let $X/k$ be a proper scheme over a field, and let
$\mathcal L$ be an invertible Module on $X$. The following conditions are
equivalent:

\noindent a. $\mathcal L$ is $\tau$-equivalent to zero;

\noindent b. For every coherent Module $F$,
\[
\chi(F\otimes\mathcal L)=\chi(F),
\]
where $\chi$ denotes the Euler--Poincaré characteristic;

\noindent b'. As in (b), but with $F=\mathcal O_Y$, where $Y$ is a closed
integral subscheme of dimension 1 in $X$;

\noindent c. For every such $Y$, writing $Y'$ for the normalized curve,
\[
\deg\mathcal L_{Y'}=0.
\]

If $X/k$ is projective and carries an ample invertible Module
$\mathcal O_X(1)$, the preceding conditions are also equivalent to:

\noindent d. For every integer $m$, $\mathcal L^{\otimes m}(1)$ is ample.

\noindent e. (If $X$ is integral.) For every pair of integers $m,n$,
\[
\chi(\mathcal L^{\otimes m}(n))=\chi(\mathcal O_X(n))
\]

\src{306}
(that is, (b) holds with $F=\mathcal L^{\otimes m}(n)$).

For the sufficiency of this last condition, observe that it says that all the
Hilbert polynomials of the $\mathcal L^{\otimes m}$ are equal. Hence, by
Mumford's criterion (iii), the $\mathcal L^{\otimes m}$ remain in a
quasi-compact subset of $\operatorname{Pic}_{X/k}$, which is (a). Conditions
(b), (b'), and (c) should be regarded as variants, on an arbitrary proper
scheme, of numerical equivalence as customarily defined on nonsingular
projective varieties. For the latter, the equivalence of (a) and (c) was of
course well known (Matsusaka).

Criterion (e) also implies:

\medskip
\noindent(ix) Let $f:X\longrightarrow S$ be a projective flat morphism whose
geometric fibers are integral. Then $\operatorname{Pic}_{X/S}^\tau$ is open
and closed in $\operatorname{Pic}_{X/S}$.

We confine ourselves to a few comments on the proof of the key results (i),
(ii), and (iii). Result (iv) is somewhat separate and is proved using only (i)
for finite surjective radicial morphisms, specifically for a Frobenius
morphism. For (i), one makes essential use of the ideas of nonflat descent
(see especially TDTE I, page 9). Since only finiteness results are sought, the
lack of effectivity criteria for descent data is harmless. Mumford has
recently proved a slightly weaker form of (iii), namely the criterion involving
all coefficients of the Hilbert polynomial. His argument is extremely simple
and is inspired by the proof of an ampleness criterion of Nakai (stated by
Nakai for nonsingular surfaces and generalized by Mumford to arbitrary
projective morphisms). It seems to me, however, that this argument is valid
only under a slight additional restriction on the fibers of $X/S$, Serre's
condition $(S_2)$, which is satisfied when the geometric fibers are normal.
This restricted criterion is then used in proving (ii): criterion (i) reduces
one to the case in which $Y/S$ is flat with geometrically integral normal
fibers, and Mumford's criterion readily reduces one to the case in which $X/S$
satisfies the same conditions. The dimension hypothesis then implies that the
geometric fibers of $Y$ and $X$ have depth $\geq2$ at their closed points,
which permits application of the “equivalence criteria” in the form given in
SGA 1962 (Exposé XII), completing the proof of (ii). Once (i) and (ii) are
available, it is not difficult to remove every normality hypothesis on the
fibers from Mumford's criterion and prove it in the stronger form stated in
(iii).

\src{307}
Finally, the proof of (i) and (ii) also shows that when $S$ is the spectrum of
a field $k$, the morphism
\[
\operatorname{Pic}_{Y/k}\longrightarrow\operatorname{Pic}_{X/k}
\]
is affine (and not merely of finite type).

\par\smallskip
\begingroup\footnotesize
\noindent\textit{The print asks that the coefficients $a_1$ and $a_2$ remain bounded, without specifying the case $n<2$. The current reading only makes the absent-coefficient convention explicit in low dimension. This is not a claim that the printed theorem is false.}\par
\endgroup
